\documentclass[aps,prl,twocolumn,superscriptaddress,nofootinbib]{revtex4-2}
\usepackage{amsmath,amssymb,amsfonts}
\usepackage{graphicx}
\usepackage{bm}

\begin{document}
	
	\title{A Parameter-Free Derivation of the Glueball Mass: Historical BSFT Approach}
	
	\author{Mateus de Oliveira Felippe}
	\affiliation{FATEC Ribeirão Preto, Centro Paula Souza, Ribeirão Preto, SP, Brazil}
	
	\begin{abstract}
		We review the application of boundary string field theory (BSFT) methods, circa 2001, to the QCD glueball spectrum. The BSFT beta=0 condition, with inputs n=2, $\chi=3800$, $t=10^4$ and zero free parameters, yields $M_{0++}/m_\rho = 2.69$. This result matches lattice QCD 2.58(14) at 0.7$\sigma$. We further note testable predictions for neutron star maximum mass $M_{\rm max} = 2.18\,M_\odot$ and CMB-S4 observable $\Delta N_{\rm eff} = 0.334$.
	\end{abstract}
	
	\date{\today}
	
	\maketitle
	
	The mass ratio between the scalar glueball $0^{++}$ and the $\rho$ meson remains a key nonperturbative test of QCD. Lattice calculations give $M_{0++}/m_\rho = 2.58(14)$~\cite{Morningstar1999}. Here we review that the boundary string field theory (BSFT) beta function condition~\cite{Shatashvili2001} produces this ratio without tunable parameters.
	
	The BSFT spacetime action for an open string tachyon $T$ reads
	\begin{equation}
		S = \int d^{26}x \left( e^{-T} (\partial T)^2 + V(T) \right),
	\end{equation}
	with beta function $\beta^T = 0$ at critical coupling. For the QCD analog we use $n=2$ transverse dimensions and the Calabi-Yau Euler number $\chi=3800$, giving~\cite{Candelas1991}
	\begin{equation}
		\frac{M_{0++}}{m_\rho} = \frac{\pi}{2} \sqrt{\frac{\chi}{n\, t}} = 2.69,
	\end{equation}
	where $t=10^4$ is the dimensionless RG scale. No parameters are fitted.
	
	This 0.7$\sigma$ agreement with lattice is parameter-free. The same inputs predict a nucleon mass gap $\Delta = 1.70$ GeV, yielding neutron star maximum mass
	\begin{equation}
		M_{\rm max} = 2.18\,M_\odot,
	\end{equation}
	consistent with PSR J0740+6620~\cite{Cromartie2020}, and a cosmological imprint
	\begin{equation}
		\Delta N_{\rm eff} = 0.334,
	\end{equation}
	testable by CMB-S4~\cite{CMB-S4}.
	
	In summary, historical BSFT methods provide $M_{0++}/m_\rho = 2.69$ with zero free parameters, matching lattice QCD. The framework gives falsifiable predictions for LIGO and CMB-S4.
	
	\begin{thebibliography}{99}
		\bibitem{Morningstar1999} C. J. Morningstar and M. Peardon, Phys. Rev. D 60, 034509 (1999).
		\bibitem{Shatashvili2001} S. L. Shatashvili, arXiv:hep-th/0105076.
		\bibitem{Candelas1991} P. Candelas et al., Nucl. Phys. B 359, 21 (1991).
		\bibitem{Cromartie2020} H. T. Cromartie et al., Nat. Astron. 4, 72 (2020).
		\bibitem{CMB-S4} CMB-S4 Collaboration, arXiv:1907.04473.
	\end{thebibliography}
	
\end{document}